% !PRINTED 174
\weberpage{174}

\begin{center}
\textbf{\S\ 53. \gesperrt{Die L\"osungen der Gleichung} $j(\omega) = j(\omega')$.}
\end{center}

Die Resultate von \S\ 51 gen\"ugen zun\"achst, um die Beziehungen der verschiedenen Werte von $\omega$ festzustellen, die zu demselben Wert von $\varkappa^2$ f\"uhren. Sei also:

\weq{1}{\frac{\vartheta_{10}^4(0,\omega')}{\vartheta_{00}^4(0,\omega')} = \frac{\vartheta_{10}^4(0,\omega)}{\vartheta_{00}^4(0,\omega)} = \varkappa^2.}

Setzen wir

\weq{2}{v = \pi\vartheta_{00}^2(0,\omega)\,u = \pi\vartheta_{00}^2(0,\omega')\,u',}

so gen\"ugen (nach \S\ 42) sowohl die Funktionen

\weq{3}{\frac{\vartheta_{00}(0,\omega)\vartheta_{11}(u,\omega)}{\vartheta_{10}(0,\omega)\vartheta_{01}(u,\omega)}, \quad \frac{\vartheta_{01}(0,\omega)\vartheta_{10}(u,\omega)}{\vartheta_{10}(0,\omega)\vartheta_{01}(u,\omega)}, \quad \frac{\vartheta_{01}(0,\omega)\vartheta_{00}(u,\omega)}{\vartheta_{00}(0,\omega)\vartheta_{01}(u,\omega)},}

als auch

\weq{4}{\frac{\vartheta_{00}(0,\omega')\vartheta_{11}(u',\omega')}{\vartheta_{10}(0,\omega')\vartheta_{01}(u',\omega')}, \quad \frac{\vartheta_{01}(0,\omega')\vartheta_{10}(u',\omega')}{\vartheta_{10}(0,\omega')\vartheta_{01}(u',\omega')}, \quad \frac{\vartheta_{01}(0,\omega')\vartheta_{00}(u',\omega')}{\vartheta_{00}(0,\omega')\vartheta_{01}(u',\omega')},}

f\"ur $x,y,z$ gesetzt, den Differentialgleichungen (1) des vorigen Paragraphen, und die entsprechenden dieser Funktionen sind also identisch. Wenn nun $u'=\frac{1}{2}$ ist, so verschwindet $\vartheta_{10}(u',\omega')$ und es mu\ss{} dann auch $\vartheta_{10}(u,\omega)$ verschwinden. Demnach folgt aus (2) mit R\"ucksicht auf \S\ 21:

\weq{5}{\vartheta_{00}^2(0,\omega') = \vartheta_{00}^2(0,\omega)(\alpha+\beta\omega),}

worin $\alpha,\beta$ ganze Zahlen sind, die der Bedingung

\weq{6}{\alpha \equiv 1, \quad \beta \equiv 0 \ (\mathrm{mod}\ 2)}

gen\"ugen. Ist $u'=\omega':2$, so verschwinden $\vartheta_{01}(u',\omega')$ und $\vartheta_{01}(u,\omega)$, und daher ist wie oben

\weq{7}{\vartheta_{00}^2(0,\omega')\,\omega' = \vartheta_{00}^2(0,\omega)(\gamma+\delta\omega),}

worin

\weq{8}{\gamma \equiv 0, \quad \delta \equiv 1 \ (\mathrm{mod}\ 2).}

Daraus folgt:

\weq{9}{\omega' = \frac{\gamma+\delta\omega}{\alpha+\beta\omega}.}

Da aber in gleicher Weise geschlossen werden kann:

\[
\begin{aligned}
\vartheta_{00}^2(0,\omega) &= \vartheta_{00}^2(0,\omega')(\alpha'+\beta'\omega'),\\
\vartheta_{00}^2(0,\omega)\,\omega &= \vartheta_{00}^2(0,\omega')(\gamma'+\delta'\omega'),
\end{aligned}
\]

worin $\alpha',\beta',\gamma',\delta'$ gleichfalls ganze Zahlen sind, so ergibt die Substitution (5) in diesen Gleichungen:

\[
\begin{aligned}
1 &= (\alpha+\beta\omega)(\alpha'+\beta'\omega'),\\
\omega &= (\alpha+\beta\omega)(\gamma'+\delta'\omega').
\end{aligned}
\]
